% ECONOMICS_PROBLEM: {"id": "D84-559", "title": "Inflation-expectation formation", "jel_code": "D84", "source_id": "OP-0243"}
% !TeX program = xelatex
\documentclass[11pt,a4paper]{article}
\usepackage[margin=23mm]{geometry}
\usepackage{fontspec}
\setmainfont{Latin Modern Roman}
\usepackage{amsmath,amssymb,mathtools}
\usepackage{xcolor,enumitem,booktabs,longtable,array,xurl,hyperref}
\definecolor{muted}{HTML}{53636C}
\hypersetup{colorlinks=true,urlcolor=blue,linkcolor=blue}
\setlength{\parindent}{0pt}
\setlength{\parskip}{5pt}
\newcommand{\R}{\mathbb R}
\newcommand{\E}{\mathbb E}
\newcommand{\N}{\mathbb N}
\newcommand{\ind}{\mathbf 1}
\newcommand{\Var}{\operatorname{Var}}
\newcommand{\Cov}{\operatorname{Cov}}
\newcommand{\supp}{\operatorname{supp}}
\DeclareMathOperator*{\argmin}{arg\,min}
\DeclareMathOperator*{\argmax}{arg\,max}
\newcommand{\entrymeta}[3]{{\sffamily\small\color{muted}\textbf{\texttt{#1}}\quad|\quad #2\par #3\par}}
\newcommand{\Problemref}[1]{\texttt{#1}}
\newcommand{\JELref}[2]{\texttt{#2} (JEL \texttt{#1})}
\begin{document}
\section*{Inflation-expectation formation}
\textbf{Problem ID:} \texttt{D84-559}\quad \textbf{JEL:} \texttt{D84}\par
% BEGIN ECONOMICS STATEMENT
\label{problem:OP-0243}
{\sffamily\small\color{muted}Original subject group: Inflation and monetary policy\par}
\label{definitions:problem:30007524}
\textbf{Audit classification:} Explicit published open question. Formation and relevance of expectations are explicit historical questions; the randomized design is editorial. Verification concerns the cited version, not first priority or proof that this exact editorial specification remains unsolved on October 8, 2026.
\subsection*{Definitions and precise research target}
\paragraph{Editorial mathematical specification.} For household \(i\), let \(m_{it}\) denote its subjective mean of next-year inflation, \(p_{it}\) its personally experienced price change, \(n_{it}\) the inflation information it receives, and \(a_{it}\) measured attention. Specify the belief-transition model
\[
m_{i,t+1}=G_\theta(m_{it},p_{it},n_{it},a_{it},h_{it})+\epsilon_{i,t+1},
\]
where \(h_{it}\) contains predetermined household characteristics, \(\theta\) is a finite-dimensional parameter, and \(\epsilon_{i,t+1}\) is a mean-zero innovation conditional on the stated information. The experimental design independently randomizes informative messages and attention incentives; personally experienced prices require a separate valid instrument or an explicit partial-identification strategy.
Identify the causal effects of information content, exposure to personally salient prices, and attention on both the mean and dispersion of subjective inflation distributions. Test whether one estimated transition rule transports to different inflation regimes and communication formats, using held-out households and occasions. Distinguish persistent updating from temporary survey-response priming and reconcile forecasts at different horizons. To establish macroeconomic relevance, connect beliefs to actual consumption or pricing through independently identified outcome equations rather than assuming that survey revisions imply behavior. The target is an empirically disciplined formation mechanism and its policy responsiveness within a declared population; unrestricted household-specific updating rules do not constitute a solution.

Each household has a subjective probability measure \(Q_{it}\) on next-year cumulative inflation; \(m_{it}=\int x\,dQ_{it}(x)\) and \(v_{it}=\int(x-m_{it})^2dQ_{it}(x)\), with finite second moments. Cross-sectional disagreement \(\operatorname{Var}_i(m_{it})\) differs from average subjective uncertainty \(E_i[v_{it}]\). Define own-price exposure by a fixed household consumption basket so that changes in basket weights do not silently alter treatment. Factorial assignment randomizes message content and attention incentives independently; an incentive is not the same as exogenously setting attention unless its only effect is through measured attention. Conditioning the innovation on information including randomized treatment identifies the transition only under stated restrictions on unobserved household effects and dynamic panel attrition.
\subsection*{Variants and assumption changes}
The following are \emph{editorial variants}, not additional published open conjectures.
\begin{enumerate}
\item \emph{Persistent learning.} Measure beliefs immediately and at several later horizons after the same message, with no repeated message treatment. Identify decay of the causal revision and distinguish learning from short-lived survey priming.
\item \emph{Distributional updating.} Model the full subjective density using a fixed parametric family, not only its mean. Estimate separately changes in disagreement and individual uncertainty and their effects on actual spending.
\end{enumerate}
\subsection*{References and source audit}
\begin{itemize}
\item Janet L. Yellen. \emph{Macroeconomic Research After the Crisis}. 2016. Locator: Section Inflation Dynamics. Primary text checked. \url{https://www.federalreserve.gov/newsevents/speech/yellen20161014a.htm}.
\item European Central Bank. \emph{Forecasting and business cycle analysis}. Undated. Locator: Understanding the inflation process, last selected area (expectation measurement and determinants). Primary text checked. \url{https://www.ecb.europa.eu/pub/economic-research/research_agenda/forecasting/html/index.en.html}.
\end{itemize}
Formation and relevance of expectations are explicit historical questions; the randomized design is editorial.
\begin{enumerate}\raggedright
\item\hypertarget{definitions:source:30007524:1}{} Janet L. Yellen. 2016. \emph{Macroeconomic Research After the Crisis}. Section Inflation Dynamics.
\url{https://www.federalreserve.gov/newsevents/speech/yellen20161014a.htm}.
\item\hypertarget{definitions:source:30007524:2}{} European Central Bank. Undated / date unverified. \emph{Forecasting and business cycle analysis}. Understanding the inflation process, last selected area (expectation measurement and determinants).
\url{https://www.ecb.europa.eu/pub/economic-research/research_agenda/forecasting/html/index.en.html}.
\end{enumerate}

\subsection*{Common conventions: Source mathematical conventions}

These conventions apply to each entry unless explicitly replaced.

\textbf{Models and identification.} A primitive model \(m\) specifies all probability laws, feasible choices, information, equilibrium and measurement rules used by its target. The admissible class \(\mathcal M\) is fixed independently of outcomes. If \(P_O(m)\) is its observable probability law and \(\tau(m)\) the target, its identified set at observable law \(P\) is
\[
 \mathcal I_\tau(P)=\{\tau(m):m\in\mathcal M,\ P_O(m)=P\}.
\]
Point identification means that this set is a singleton. An empty set signals incompatibility of data and assumptions. Coordinate bounds need not describe the joint set. Bounds are sharp only if every retained target is attainable or arbitrarily approachable by compatible models. Finite bounds require explicit restrictions on unobserved quantities; unrestricted confounding, hidden wealth, extrapolation or arbitrary equilibrium selection can yield uninformative sets.

\textbf{Causal interventions.} A potential outcome \(Y_i(p)\) is the outcome under a complete policy regime \(p\), including timing, financing, information and market spillovers. The target population and units are fixed before treatment. With interference, use the complete assignment vector or a justified exposure mapping; individual no-interference notation alone cannot identify an economy-wide intervention. A difference of mean potential outcomes does not require the cross-world joint distribution. Conditioning on post-treatment migration, survival, enrollment or employment changes the estimand and needs separate justification.

\textbf{Statistical statements.} An estimator, confidence set, forecast or test specifies a sampling law, model class and loss/coverage criterion. Uniform \((1-\alpha)\)-coverage means \(\inf_{m\in\mathcal M}\Pr_m\{\tau(m)\in C_n\}\geq1-\alpha\), or an explicitly asymptotic version. The whole parameter space trivially has coverage, so informativeness requires a stated width, risk or power objective. A forecasting improvement is relative to a named benchmark, information set and evaluation population.

\textbf{Equilibrium and optimization.} An equilibrium comprises feasible plans, optimality, beliefs where needed, budget constraints and all market/resource clearing. The primitives or a stated benchmark define each correspondence \(\mathcal E(p;m)\). Nonemptiness is assumed or proved, never inferred just from compact policy sets. A selected-equilibrium target uses a declared selection rule; a robust target quantifies over all permitted equilibria. Use suprema or infima unless attainment is established. An \(\varepsilon\)-optimal policy is feasible and within \(\varepsilon\) of the finite optimum value. Report an empty feasible set separately.

\textbf{Welfare and accounting.} Normative utility, interpersonal normalization, social weights, numeraire, discounting and the use of public receipts are primitives. Behavioral utility need not coincide with the welfare criterion. Household welfare includes ownership income and rents once; adding profits again to compensating variation generally double-counts them. A monetary equivalent is defined by an explicit transfer and price convention, with monotonicity and range conditions. Average effects, mechanism contributions, transfers and welfare losses are different targets. Infinite sums require convergence and dynamic feasible plans require terminal or no-Ponzi conditions.

\textbf{Source versions.} A working-paper date, revision date and journal publication year are distinguished. A DOI for a journal version is not automatically the DOI of an earlier manuscript. Locators refer to the stated version and printed pages where specified. A metadata-only check does not verify a claim about a full-text conclusion. Every entry's source subsection narrows the provenance claim accordingly.
% END ECONOMICS STATEMENT
\end{document}
